\documentclass[12pt]{article}
\usepackage[utf8]{inputenc}
\usepackage[T1]{fontenc}
\usepackage{hyperref}
\title{Itang Language Resources for K-12 Curriculum Development}
\author{Compiled by Claude Sonnet 5 (claude-sonnet-5) for Andr\'as Kornai}
\date{}
\begin{document}
\maketitle
\section{Introduction}
Itang (ISO 639-3 \texttt{uta}) is also known as Otank, Itank, Otang and Utanga. Wikipedia classes it as a Tivoid language of Nigeria. It is also reported in Cameroon. Speaker estimates range from 15,000 (Wikipedia, 2006) to about 25,000 (Joshua Project). A 2023 conference presentation cites 24,000 speakers and calls the language under-described. Wikipedia states that it has no writing system. Very little material exists, and none of what I found was written for schools.
\section{Primary Educational Resources}
I found no primers, textbooks or school programmes for Itang.
\section{Linguistic and Reference Materials}
\begin{itemize}
\item Otank language (Wikipedia overview): \url{https://en.wikipedia.org/wiki/Otank_language}
\item Itang noun class system (Graz, 8 Dec 2023), a conference presentation. It says no dedicated linguistic study of Itang had been done as of then, and names two dialects, Kundeve and Besenge: \url{https://nuclaba.blogs.uni-hamburg.de/wp-content/uploads/2025/10/Itang-noun-class-system-Graz-20231208.pdf}
\end{itemize}
\section{Religious and Community Resources}
\begin{itemize}
\item Scripture Earth entry for code uta, which lists Itang among the alternative names: \url{https://www.scriptureearth.org/00eng.php?iso=uta}
\item Joshua Project language page: \url{https://www.joshuaproject.net/index.php/languages/uta}
\end{itemize}
I did not check what scripture material, if any, exists behind these listings.
\section{Digital and Audio Resources}
I found no audio, video or corpus resources in the search results.
\section{Curriculum Development Context}
The language has no established orthography and no pedagogical materials that I could find. Any K-12 curriculum would first need community-led orthography development and documentation. The Graz study on noun classes is a starting point for linguists.
\end{document}